% Artigo: O Câmbio na Imprensa — posicionamentos dos jornais brasileiros
% sobre a Caixa de Conversão (1906-1914)
% Pedro Ortencio Pires de Campos Telles — FFLCH-USP
%
% Compilar: pdflatex main && bibtex main && pdflatex main && pdflatex main

\documentclass[12pt,a4paper]{article}

\usepackage[utf8]{inputenc}
\usepackage[T1]{fontenc}
\usepackage[brazil]{babel}
\usepackage{graphicx}
\usepackage{booktabs}
\usepackage{natbib}
\usepackage[margin=2.5cm]{geometry}
\usepackage{setspace}
\onehalfspacing

\title{O Câmbio na Imprensa: os posicionamentos dos jornais brasileiros\\
sobre a Caixa de Conversão (1906--1914)}
\author{Pedro Ortencio Pires de Campos Telles\thanks{FFLCH-USP, História Econômica. pedrortencio@gmail.com}}
\date{\today}

\begin{document}
\maketitle

\begin{abstract}
% Resumo: pergunta (qual pensamento prevaleceu no debate da Caixa entre os
% principais diários de RJ/SP), corpus (4 jornais, Hemeroteca Digital, 1906-1914),
% método (stance detection por LLM validada por codificação humana, DSL),
% principais resultados. 150-250 palavras.
\end{abstract}

\section{Introdução}\label{sec:intro}
% Seção: introducao — a redigir (skill escrita-academica; sem travessões)


\section{O debate monetário e a Caixa de Conversão, 1906--1914}\label{sec:contexto}
% Seção: contexto — a redigir (skill escrita-academica; sem travessões)


\section{Dados e método}\label{sec:metodo}
% Seção: metodo — a redigir (skill escrita-academica; sem travessões)


\section{Resultados}\label{sec:resultados}
% Seção: resultados — a redigir (skill escrita-academica; sem travessões)


\section{Conclusão}\label{sec:conclusao}
% Seção: conclusao — a redigir (skill escrita-academica; sem travessões)


\bibliographystyle{apalike}
\bibliography{referencias}

\end{document}
